(1) 令作用在球1上的冲量为 \(+J\mathbf n\)，作用在球2上的为 \(-J\mathbf n\)。冲量动量定理给 \(\mathbf v_1'=\mathbf v_1+(J/m_1)\mathbf n\)、\(\mathbf v_2'=\mathbf v_2-(J/m_2)\mathbf n\)。恢复系数要求 \((\mathbf v_1'-\mathbf v_2')\cdot\mathbf n=-e g_n\)，代入后 \(g_n+J(1/m_1+1/m_2)=-e g_n\)，故 \(J=-(1+e)\mu g_n>0\)。它是法向排斥冲量，与 \(g_n<0\) 的碰撞前接近约定相符。

(2) 相对速度差为 \(\mathbf g'=\mathbf g+(J/\mu)\mathbf n=\mathbf g-(1+e)g_n\mathbf n\)。因此分解成 \(\mathbf g_\parallel=g_n\mathbf n\) 和 \(\mathbf g_\perp=\mathbf g-g_n\mathbf n\) 后有 \(\mathbf g_\perp'=\mathbf g_\perp\)、\(g_n'=-eg_n\)。总动量守恒使质心动能不变，损失只来自相对法向动能，故 \(K_{\rm before}-K_{\rm after}=\frac12\mu(g_n^2-e^2g_n^2)=\mu(1-e^2)g_n^2/2\ge0\)。若有人写出切向动能损失，则与光滑接触无摩擦条件冲突。

(3) 质心速度 \(\mathbf V=(m_1\mathbf v_1+m_2\mathbf v_2)/(m_1+m_2)\) 在冲击前后相同。总动能可严格分解为 \(\frac12(m_1+m_2)V^2+\frac12\mu|\mathbf g|^2\)。弹性碰撞 \(e=1\) 时法向相对速度反号而模长不变，因此相对动能完全保持；完全非弹性法向 \(e=0\) 时两球法向分速度相同，但切向仍可不同，不能自动断言它们黏合或全部相对速度为零。

(4) 令球2碰前静止，球1速度 \(\mathbf v_1=v(\cos\beta\,\mathbf n+\sin\beta\,\mathbf t)\)，其实际接近要求 \(\cos\beta<0\)。则 \(g_n=v\cos\beta\)，冲量 \(J=-(1+e)\mu v\cos\beta\)，球1碰后法向分量 \(v'_{1n}=v\cos\beta[1-(1+e)m_2/(m_1+m_2)]\)，切向分量 \(v'_{1t}=v\sin\beta\)。若取一维正碰 \(\beta=\pi\)，两个球的法向速度遵从标准一维恢复系数公式；\(e=1,m_1=m_2\) 时入射球法向分量消失并由静止球接收，和等质量弹性正碰一致。
